% ==========================================================================
%  Chapter 99 -- Branch Selection at Softening Reversals (Extension Plan, §F0)
% ==========================================================================

\chapter{Branch Selection at Softening Reversals: Problem, Theory and Plan}
\label{ch:branch-selection-reversals}

\paragraph{Normative status.}
This chapter opens the \emph{branch-selection extension} of the
regularized continuation solver. It is a research growth path, not a
correction of a defect: the \(\pm200\,\mathrm{mm}\) cyclic protocol on
the reduced RC column already \emph{closes} under Levenberg--Marquardt
(LM) continuation (Cap.~\ref{ch:arc-length} and the validation
chapters), but a fraction of the reversal steps are accepted as
\emph{regularized} rather than fully converged, and the accepted branch
at some reversals carries a spurious force spike. The extension adds
three independent, \textbf{env-gated, default-off} remedies that
\emph{select} the physical branch: deflated continuation
(Cap.~\ref{ch:deflated-continuation}), energy-minimisation continuation
(Cap.~\ref{ch:energy-tao-continuation}), and a Cultural-Algorithm
meta-tuner built on a new std-only optimisation module
(Cap.~\ref{ch:cultural-algorithm-module}). Every remedy is compiled out
of the default build path, so with all gates off the solver is
bit-identical to the baseline recorded here.

\section{The baseline and what it exposes (Gate G0)}

The reference run \texttt{g0\_baseline} reproduces the current solver on
a \(2\times2\times4\) Hex27 mesh, \texttt{profile=production\_stabilized},
\texttt{tangent=fracture\_secant}, over the four-amplitude cyclic
protocol \(\{50,100,150,200\}\,\mathrm{mm}\):
%
\begin{center}
\begin{tabular}{@{}ll@{}}
\toprule
Accepted steps            & \(196\) \\
Regularized (not fully converged) & \(80\) \\
Crosses the \(88\,\mathrm{mm}\) limit point & yes \\
Peak \(|V_\text{base}|\)  & \(5.08\times10^{-2}\,\mathrm{MN}\) \\
\bottomrule
\end{tabular}
\end{center}
%
The \(80/196\) figure is the quantitative target of the extension: a
successful branch-selection remedy must lower the count of regularized
reversal steps (or the reversal spike magnitude) \emph{without}
disturbing the well-conditioned unloading segments. The artefacts live
under \path{data/output/g0_baseline/cyc/} and the driver is
\href{run:src/validation/ReducedRCColumnContinuumBaseline.cpp}{%
\texttt{ReducedRCColumnContinuumBaseline.cpp}}.

\section{Why a reversal admits several equilibria}

Let the discrete equilibrium at a frozen material state be
\(\mathbf{R}(\mathbf{u})=\mathbf{f}_\text{int}(\mathbf{u})
-\lambda\,\mathbf{f}_\text{ext}=\mathbf{0}\). For a hardening material
\(\mathbf{R}\) is (locally) strictly monotone and the root is unique. The
Ko--Bathe continuum is different: it combines strain \emph{softening}
with a \emph{fixed} smeared crack and unilateral (tanh-smoothed) closure
(Cap.~\ref{ch:ko-bathe-3d}). Once a Gauss point has cracked, its
tangent contribution can be negative, and the crack-band regularisation
fixes the localisation width but not the number of admissible global
states. As a consequence the softening branch of \(\mathbf{R}=\mathbf{0}\)
is \emph{non-unique}: the classic result that smeared softening produces
mesh-sensitive, multiply-valued equilibrium paths is exactly the
crack-band and rotating/fixed-crack literature
\cite{BazantOh1983CrackBand,RotsBlaauwendraad1989CrackModels}.

At a load reversal the situation sharpens. The warm-start
\(\mathbf{u}_k\) sits at the softened extreme of the previous half-cycle.
Two equilibria are then reachable from it at the new control value:
%
\begin{enumerate}
  \item the \emph{physical} low-force branch, where the previously open
        cracks unload elastically and the section recovers normal
        stiffness on closure;
  \item a \emph{spurious} high-force branch, where the iteration relaxes
        into a different combination of open/closed cracks that also
        satisfies \(\mathbf{R}=\mathbf{0}\) but at a much larger base
        shear.
\end{enumerate}
%
Both are stationary points of the incremental potential
\(\Pi(\mathbf{u})\) whose gradient is \(\mathbf{R}=\nabla\Pi\) (with the
state frozen inside the step, \(\mathbf{R}\) is a gradient because the
secant modulus is symmetric, Cap.~\ref{ch:ko-bathe-3d}). The
physical branch is a \emph{minimiser} of \(\Pi\); the spurious high-force
branch is typically a saddle or a maximiser along the reversal
direction. This distinction is the lever the energy remedy pulls
(Cap.~\ref{ch:energy-tao-continuation}).

\section{Why regularization alone does not select the branch}

The LM continuation \cite{Levenberg1944,Marquardt1963} solves
\((\mathbf{K}+\mu\mathbf{I})\,\delta\mathbf{u}=-\mathbf{R}\) with an
adaptive \(\mu\). Its purpose is \emph{conditioning}, not selection: a
positive \(\mu\) makes a near-singular tangent invertible so the
continuation crosses the limit point where pure Newton stalls, and
\(\mu\!\to\!0\) on a well-conditioned branch recovers a local Newton
step. Crucially, LM converges into whatever basin of attraction the
warm-start lies in. At a reversal the plain warm-start \(\mathbf{u}_k\)
can lie in the basin of the spurious branch, and no amount of tangent
regularisation moves it out: \(\mu\) only rescales the step, it does not
change which root the iteration is attracted to. The secant predictor
already implemented (Cap.~\ref{ch:arc-length},
\href{run:src/analysis/RegularizedNewtonContinuation.hh}{%
\texttt{RegularizedNewtonContinuation.hh}}) mitigates this by aiming the
warm-start down the unloading branch, but it is a heuristic and does not
\emph{guarantee} branch selection.

Branch selection therefore requires a mechanism that is \emph{aware of
the other roots}. Three such mechanisms are developed.

\section{The three remedies}

\paragraph{(P1) Deflated continuation.}
Deflation multiplies the residual (equivalently, rescales the Newton
direction) by an operator that becomes singular at previously found
roots, so the iteration is \emph{repelled} from any root registered as
spurious \cite{FarrellBirkissonFunke2015}. Registering the spurious
high-force reversal state and re-solving forces convergence to a
different --- physical --- root. Developed and verified in isolation in
Cap.~\ref{ch:deflated-continuation}.

\paragraph{(P2) Energy-minimisation continuation.}
Instead of hunting a zero of \(\mathbf{R}\), minimise the incremental
potential \(\Pi\) with a trust-region Newton optimiser
\cite{NocedalWright2006,ConnGouldToint2000} from PETSc/TAO
\cite{BalayPETScManual}. A descent method cannot terminate at a maximiser
or saddle, so it selects the physical minimiser branch by construction.
Developed in Cap.~\ref{ch:energy-tao-continuation}.

\paragraph{(M/P3) Cultural-Algorithm meta-tuning.}
The solver exposes a small parameter vector (initial \(\mu\), growth and
drop factors, stagnation window, predictor scale, acceptance floor,
reversal subdivision). A Cultural Algorithm \cite{Reynolds1994Cultural,
ReynoldsPeng2004} searches this space by replaying a short window across
a reversal, scoring by converged-fraction and spike magnitude. The
algorithm is implemented as a reusable, dependency-free optimisation
module (Cap.~\ref{ch:cultural-algorithm-module}) that will also serve the
control-device design problem of
\cite{LaraValenciaEchavarriaValencia2024}.

\section{Design invariants (binding)}

\begin{enumerate}
  \item \textbf{Default equivalence.} Every remedy sits behind a
        compile-time gate (\texttt{if constexpr} on a policy flag) or an
        environment flag defaulting to off. With all gates off the CSV
        output is bit-identical to \texttt{g0\_baseline}.
  \item \textbf{No physics rewrite.} The remedies drive the existing
        assemble/solve/accept seam; neither the constitutive law nor the
        element nor PETSc SNES is re-implemented.
  \item \textbf{Cheap rollback.} A spurious step is detected only
        \emph{after} it is committed (the crack ratchet has advanced), so
        every retry is bracketed by a state checkpoint/restore.
  \item \textbf{Isolated verifiability.} Each remedy ships a std-only or
        \texttt{CORE}-tier unit test that exercises the mechanism on a toy
        problem before any coupling to the full column.
\end{enumerate}

These invariants are what make the extension safe to carry in the
mainline solver header without endangering the validated baseline.
